\section{Introduction}
\label{sec:intro}

How much synthetic data is enough?
Model-generated data is now a default way to stretch a small real dataset, whether for distillation, augmentation, or instruction tuning.
The cost of generation grows linearly with the amount generated, so the choice of how much to generate is a budget decision.
If the benefit saturates early, most of the generation budget is wasted.
If large amounts actively hurt, the budget buys harm.

Two lines of prior work bear on this question, and neither answers it.
On one hand, work on model collapse shows that training generative models recursively on their own outputs loses the tails of the data distribution~\citep{shumailov2023curse,shumailov2024collapse,alemohammad2024mad,dohmatob2024tale}, and that keeping real data in the mix limits the damage~\citep{gerstgrasser2024inevitable,bertrand2024stability}.
These studies concern many generations of self-training.
The common practical setting is different: a single round of generation by a \emph{stronger} model, added to a fixed real set.
On the other hand, synthetic augmentation studies report scaling curves for image recognition~\citep{he2023synthetic,azizi2023synthetic,fan2024scaling} and text classification~\citep{ye2022zerogen,yu2023attributed,li2023synthetic}, but each covers one task family with its own generator, student, and metric.
Their curves cannot be compared across domains.

We run one protocol across three task families (\figref{fig:pipeline}).
A teacher model sees three random examples from a 100-example real seed set and writes one new example per request.
We fine-tune a smaller student on the 100 real examples plus $r \times 100$ synthetic ones, for $r$ from 0 to 50, and evaluate on \emph{real} held-out data.
We cover code generation (\mbpp), math word problems (\gsm), and five-sentence story writing (\roc), with three seeds, an all-real control at matched training-set size, and a direct measurement of the diversity of each synthetic pool.

We use this protocol to test a specific hypothesis in three parts.
(H1) Gains from synthetic data are logarithmic in the ratio, and a \ratio{3} ratio captures at least 80\% of the maximum benefit.
(H2) Beyond \ratio{10}, performance plateaus or degrades because the synthetic distribution collapses onto a few modes.
(H3) Code generation is the most resilient to high ratios, and creative writing degrades earliest.

{\bf Is there a single ceiling?} No.
At \ratio{3} the student has 63\% of its maximum observed gain on code (95\% CI 13--100\%) and 40\% on math (CI 6--64\%, which excludes 80\%).
Beyond \ratio{10} the three domains part ways.
Code pass@1 is flat ($+0.9$ points from \ratio{10} to \ratio{30}, CI $-1.6$ to $+3.6$).
Math accuracy rises from 38.3\% at \ratio{10} to 41.7\% at \ratio{50}.
Fit to human stories gets worse at every ratio we tested: 100 extra real stories lower the loss on human stories by 0.017 nats per token, while 5,000 synthetic stories raise it by 0.101.
On a metric shared by all three tasks, the loss on real held-out references, every domain shows an interior optimum followed by degradation past \ratio{10}.

{\bf Is mode collapse real in a single round of generation?} Yes, and it is our most robust finding.
In a 3,000-example synthetic code pool, 44\% of task descriptions are exact duplicates and 79\% have a near-duplicate.
In one math pool, 24\% of problems begin with the name ``Sarah''.
The effective number of distinct examples in a synthetic pool saturates at roughly half the level of real data.
Students inherit these modes: ``Sarah'' appears 171 times in 500 stories from a student trained at \ratio{30}, against 11 times after real-only training.

In summary, our contributions are as follows:
\begin{itemize}[leftmargin=*,itemsep=0pt,topsep=0pt]
    \item We conduct a controlled dose--response study of the synthetic-to-real ratio (0 to \ratio{50}) across code, math, and creative writing with one teacher, one student, one real-set size, and an all-real control at matched size (117 fine-tuning runs).
    \item We show that the ``80\% at \ratio{3}'' rule fails in every domain where it is defined, and that the behavior beyond \ratio{10} depends on both the domain and the metric: code accuracy plateaus, math accuracy keeps rising, and fidelity to real references degrades everywhere.
    \item We measure generator mode collapse directly and trace it into student outputs, including a homogenization of story content that distinct-$n$ statistics miss.
    \item We report an unexpected result: at equal training-set size, the \ratio{30} synthetic mix beats 3,100 real \gsm examples by 6.5 accuracy points (CI $+2.3$ to $+11.0$).
\end{itemize}

We review related work in \secref{sec:related}, describe the protocol in \secref{sec:method}, present results in \secref{sec:results}, and discuss implications and limitations in \secref{sec:discussion}.

\begin{figure}[t]
    \centering
    \begin{tikzpicture}[
        box/.style={draw, rounded corners=2pt, align=center, font=\footnotesize, inner sep=3pt, minimum height=0.9cm},
        arr/.style={-{Stealth[length=2mm]}, thick},
        node distance=0.5cm
    ]
        \node[box, fill=humanpurple] (real) {100 real\\examples};
        \node[box, right=of real] (teacher) {Teacher (3B)\\3-shot, $T{=}1$};
        \node[box, fill=aigreen, right=of teacher] (pool) {Synthetic pool\\(no de-duplication)};
        \node[box, right=1.1cm of pool] (student) {Student (0.5B)\\LoRA fine-tune};
        \node[box, fill=humanpurple, right=of student] (eval) {Real held-out\\test data};
        \draw[arr] (real) -- (teacher);
        \draw[arr] (teacher) -- (pool);
        \draw[arr] (pool) -- node[above, font=\scriptsize] {$r{\times}100$} (student);
        \draw[arr] (student) -- (eval);
        \draw[arr] (real.south) -- ++(0,-0.45) -| node[pos=0.25, below, font=\scriptsize] {100 real examples, always included} (student.south);
    \end{tikzpicture}
    \caption{Protocol. A teacher sees three random real examples per request and writes one new example. We fine-tune a student on 100 real examples plus $r \times 100$ synthetic ones, for $r \in \{0, 1, 2, 3, 5, 10, 20, 30, 50\}$, and evaluate on real held-out data. We repeat this for code, math, and story writing with three seeds, each with its own real seed set and its own synthetic pool.}
    \label{fig:pipeline}
\end{figure}
